\subsection*{NT-E06. Uniform CM angle counts from one Hecke-power family}
Let $E/k$ be an elliptic curve whose complex multiplication is defined
over the number field $k$. Put
\[
 a_{\mathfrak p}=N\mathfrak p+1-\#E_{\mathfrak p}(k_{\mathfrak p}),
 \qquad t_{\mathfrak p}=\frac{a_{\mathfrak p}}{\sqrt{N\mathfrak p}}.
\]
Only primes of good reduction are counted. The classical CM character
theorem supplies a unitary Hecke character $\eta$ of $k$ for which
\[
 t_{\mathfrak p}=\eta(\mathfrak p)+\overline{\eta(\mathfrak p)}
                =2\cos(2\pi\theta_{\mathfrak p}),
 \qquad \theta_{\mathfrak p}\in\mathbb R/\mathbb Z.       \tag{NT-E06a}
\]
Assume RH for $L_k(s,\eta^m)$ for every integer $m$, including
$m=0$, where the function is $\zeta_k$. This is a precisely specified
subfamily of $M$; Dirichlet GRH alone does not supply its nonzero
infinity types. Then, uniformly over all arcs $J\subset\mathbb R/\mathbb Z$
and all intervals $I\subset[-2,2]$, for $x\ge2$,
\begin{align}
 \#\{\mathfrak p:N\mathfrak p\le x,\ \theta_{\mathfrak p}\in J\}
   &= |J|\operatorname{Li}_2(x)
       +O_{E,k}\bigl(\sqrt x\log^2(2x)\bigr), \tag{NT-E06b}\\
 \#\{\mathfrak p:N\mathfrak p\le x,\ t_{\mathfrak p}\in I\}
   &= \mu(I)\operatorname{Li}_2(x)
       +O_{E,k}\bigl(\sqrt x\log^2(2x)\bigr), \tag{NT-E06c}\\
 d\mu(t)&=\frac{dt}{\pi\sqrt{4-t^2}}.\notag
\end{align}
The constants are independent of the interval and its endpoint choices.
In particular the intervals may vary with $x$. No numerical implicit
constant or uniformity in the curve or base field is asserted.

\paragraph{Classical input and complete deduction.}
We use the Hecke prime-sum estimate recorded by Bucur--Fit\'e--Kedlaya,
\emph{Effective Sato--Tate conjecture for abelian varieties and
applications}, Proposition2.5, equation(2.3), published p.1720
(author v2 p.6). Its CM interpretation is explicitly identified in
footnote2 to Proposition4.7, published p.1740 (author p.20).
For the fixed curve and field it gives, uniformly in $m\ne0$ and $x\ge2$,
\begin{align}
 S_m(x):=\sum_{N\mathfrak p\le x}\eta(\mathfrak p)^m
       &=O_{E,k}\bigl(\sqrt x\log(2x+2|m|)\bigr),\notag\\
 S_0(x)&=\operatorname{Li}_2(x)
       +O_{E,k}\bigl(\sqrt x\log(2x)\bigr).               \tag{NT-E06d}
\end{align}
Their sums omit the primes over the absolute conductor; restoring any
good primes omitted by that convention changes a sum by $O_{E,k}(1)$.
The representation dimension in this specialization is one, while the
motivic weight is $|m|$: the distinction is the point of their
Remark2.7. The finite conductor of $\eta^m$ divides that of $\eta$,
and its archimedean parameters grow at most linearly in $|m|$.
Every nonzero power has a nonzero angular infinity type, so is not
a pure norm twist. Consequently the required entire continuation,
functional equation, and boundary nonvanishing are classical Hecke
facts; the only conjectural analytic input in (NT-E06d) is the
stated RH family. The standard Hecke completion over $k$ has
degree $[k:\mathbb Q]$, which is fixed here.

For clarity, the elementary harmonic-analysis input is
Erd\H{o}s--Tur\'an, \emph{On a problem in the theory of uniform
distribution.I} (1948), TheoremIII, printed p.374 (parallel
proceedings pagination p.1150). For $P$ points on the circle and an
integer $H\ge1$ its discrepancy bound is
\[
 \left|\#\{\theta_j\in J\}-|J|P\right|
 \le C\left(\frac P{H+1}
       +\sum_{m=1}^{H}\frac1m
                    \left|\sum_{j=1}^{P}e^{2\pi i m\theta_j}\right|\right),
                                                               \tag{NT-E06e}
\]
where $C$ is absolute. The original statement includes closed
endpoints. The same bound, after changing $C$, holds with any endpoint
convention: shrink or enlarge an arc, use the same bound, and pass to
the limit for the finite set of points. Singletons are handled by this
argument as well.

Apply (NT-E06e) with $H=\lceil x\rceil$ and the angles in (NT-E06a).
The first term is $O_k(1)$, since the number of prime ideals of norm at
most $x$ is at most $[k:\mathbb Q]x$. By (NT-E06d) the sum is at most
\[
 C_{E,k}\sqrt x\log(4x+2)\sum_{m\le\lceil x\rceil}\frac1m
     =O_{E,k}\bigl(\sqrt x\log^2(2x)\bigr).
\]
Replacing $P$ by $\operatorname{Li}_2(x)$ using $m=0$ proves
(NT-E06b). The inverse image of a trace interval under
$\theta\mapsto2\cos(2\pi\theta)$ is the union of at most two arcs.
The pushforward of normalized Lebesgue measure is exactly $\mu$;
adding the two estimates proves (NT-E06c), with singleton endpoints
already accounted for. This proves all the stated uniformities.

\paragraph{Geometric meaning and locator.}
The local system $R^1f_*\mathbb Q_\ell$ of a smooth elliptic scheme over
an arithmetic open has rank two and weight one. CM splits its scalar
extension into two conjugate rank-one characters; their unitary
normalizations give the circle above. Its tensor powers are the
specific Hecke family used in the proof. Thus this is a Frobenius
distribution statement for a CM torus in the motivic monodromy group,
not a deduction of GRH for an arbitrary non-CM elliptic $L$-function.
Work/run: \texttt{GRH\_NIGHT/number\_theory},
\texttt{cm\_prime\_angles.tex}, NT-E06; sources S-NT-BFK and S-NT-ET.

\subsection*{NT-E07. A least-prime bound in any positive-measure CM interval}
Under the exact hypotheses of NT-E06, for any trace interval $I$ with
$\mu=\mu(I)>0$, there is a good prime ideal with $t_{\mathfrak p}\in I$
and
\[
 N\mathfrak p\le C_{E,k}\,\mu^{-2}\log^6(e/\mu).           \tag{NT-E07}
\]
The analogous conclusion holds for an arc of length $\mu$. This is a
newly proved formulation of the classical discrepancy consequence;
no claim of literature novelty or numerical value of $C_{E,k}$ is made.

\paragraph{Complete proof.}
Let $A_{E,k}$ be a common error constant in (NT-E06b)--(NT-E06c).
For $x\ge4$,
$\operatorname{Li}_2(x)\ge x/(2\log x)$, by integrating over $[x/2,x]$.
Write $L=\log(e/\mu)\ge1$ and set $x=C\mu^{-2}L^6$, with $C\ge4$.
Since $\log L\le L$ and $\log(1/\mu)\le L$, we have
\[
 \log(2x)\le(\log(2C)+8)L,\qquad
 \mu\sqrt x=\sqrt C\,L^3.
\]
The prime count is positive provided
$\mu\sqrt x>2A_{E,k}\log x\,\log^2(2x)$.
It is therefore enough to choose the fixed $C$ so that
\[
 \sqrt C>2A_{E,k}\bigl(\log(2C)+8\bigr)^3.
\]
Such a $C$ exists because $\sqrt C/\log^3(2C)\to\infty$.
This choice is independent of $I$, proves the bound, and also covers
intervals with endpoints included or excluded. Geometrically it
detects a prescribed positive-measure region in the compact CM
monodromy torus by one closed arithmetic fibre. Work/run:
\texttt{GRH\_NIGHT/number\_theory}, \texttt{cm\_prime\_angles.tex}, NT-E07.

\subsection*{NT-E08. An extremal CM trace asymptotic with a power-saving error}
Let now $E/k$ be any elliptic curve with potential CM by the imaginary
quadratic field $K$. Put $F=kK$ and $d=[F:k]\in\{1,2\}$. Let $\eta$
be a CM unitary Hecke character of $E_F$. Assume RH for the functions
$L_F(s,\eta^m)$ for every integer $m$, including $m=0$.
For either choice of sign, define
\[
 M^\pm_{E/k}(x)
  =\#\{\mathfrak p\text{ good}:N\mathfrak p\le x,\
               a_{\mathfrak p}=\pm\lfloor2\sqrt{N\mathfrak p}\rfloor\}.
\]
Then
\begin{align}
 M^\pm_{E/k}(x)
   &=\frac1{\pi d}\int_2^x\frac{t^{-1/4}}{\log t}\,dt
       +O_{E,k}\bigl(x^{5/8}\log^2(2x)\bigr),             \tag{NT-E08a}\\
 M^\pm_{E/k}(x)
   &\sim \frac4{3\pi d}\frac{x^{3/4}}{\log x}.
                                                               \tag{NT-E08b}
\end{align}
In particular $M$ supplies the result conditionally. The factor is
$4/(3\pi)$ when all CM is defined over $k$, and $2/(3\pi)$ in the
quadratic descent case. The integral in (NT-E08a) is retained because
replacing it by its first asymptotic term would introduce a larger
logarithmic error. This is a fully proved formulation in this run,
not an assertion that the asymptotic or its constant is new in the
literature.

\paragraph{Complete proof when $d=1$.}
It suffices first to consider the positive sign. If $a$ is an integer
with $a\le y$, then $a=\lfloor y\rfloor$ if and only if $a>y-1$.
Thus the Hasse bound gives the exact identity
\[
 a_{\mathfrak p}=\lfloor2\sqrt{N\mathfrak p}\rfloor
 \quad\Longleftrightarrow\quad
 t_{\mathfrak p}\in I_{N\mathfrak p},\qquad
 I_u=(2-u^{-1/2},2].                                    \tag{NT-E08c}
\]
The strict lower endpoint is essential and makes the identity valid
also when $N\mathfrak p$ is a square. No fractional-part factor occurs.
Direct integration of the arcsine density gives, uniformly for $u\ge2$,
\[
 \mu(I_u)=\frac1\pi\arccos\!\left(1-\frac1{2\sqrt u}\right)
         =\frac1\pi u^{-1/4}+O(u^{-3/4}).                 \tag{NT-E08d}
\]
For example the last estimate follows by setting $v=u^{-1/2}$ and
integrating $(4z-z^2)^{-1/2}=(2\sqrt z)^{-1}(1+O(z))$
from $0$ to $v$; the constants are absolute on $0\le v\le2^{-1/2}$.

For $x\ge4$ set $h=x^{-1/8}$ and partition $[2,x]$ into intervals
$(u_{j+1},u_j]$, with $u_0=x$, successive ratios $u_j/u_{j+1}=1+h$,
and the final lower endpoint truncated to $2$. Every ratio is at most
$1+h$. Since $I_u$ decreases with $u$, the contribution on this
interval is bracketed by the counts using the two fixed intervals
$I_{u_j}$ and $I_{u_{j+1}}$. Apply (NT-E06c) at both endpoints.
After summing, the total error in each bracket is at most
\[
 C_{E,k}\sum_j\sqrt{u_j}\log^2(2u_j)
       \ll_{E,k}h^{-1}\sqrt x\log^2(2x),                 \tag{NT-E08e}
\]
because $\sum_{j\ge0}(1+h)^{-j/2}\ll h^{-1}$ for $0<h\le1$.
Primes of norm exactly $2$ contribute $O_k(1)$.

The main brackets are the upper and lower Riemann sums of the
decreasing function $\mu(I_u)$ against the positive measure
$d\operatorname{Li}_2(u)=du/\log u$. Replacing $\mu(I_u)$ using
(NT-E08d) introduces an error bounded by a constant times
\[
 \int_2^x\frac{u^{-3/4}}{\log u}\,du+1
      \ll x^{1/4}+1.                                   \tag{NT-E08f}
\]
Indeed values of $u^{-3/4}$ at the two endpoints differ by a factor
at most $(1+h)^{3/4}$. For $f(u)=u^{-1/4}$ the corresponding
upper and lower sums differ by at most
$((1+h)^{1/4}-1)$ times the lower sum. They therefore bracket
$\int_2^x u^{-1/4}\,du/\log u$ with error
\[
 \ll h\int_2^x\frac{u^{-1/4}}{\log u}\,du
 \ll h x^{3/4}.                                       \tag{NT-E08g}
\]
Combining (NT-E08e)--(NT-E08g), then substituting $h=x^{-1/8}$,
proves (NT-E08a) when $d=1$. Applying exactly the same argument to
$-t_{\mathfrak p}$ proves the negative sign, since $\mu$ is symmetric.
Finally l'H\^opital's rule, or a single integration by parts, gives
\[
 \int_2^x\frac{u^{-1/4}}{\log u}\,du
      \sim \frac43\frac{x^{3/4}}{\log x},
\]
and $x^{5/8}\log^2(2x)=o(x^{3/4}/\log x)$.

\paragraph{Complete quadratic descent.}
The field defining all endomorphisms is $F=kK$. One can see this
directly on invariant differentials: in characteristic zero the
differential embeds $\operatorname{End}(E_{\bar k})\otimes\mathbb Q=K$
faithfully in $\bar k$, compatibly with $G_k$. Its kernel of Galois
action is therefore $G_{kK}$. If $d=2$, at an unramified good prime
inert in $F/k$, Frobenius interchanges the two CM eigenspaces of the
Tate module. Its matrix has zero diagonal, so $a_{\mathfrak p}=0$.
Such a prime is never extremal because $\lfloor2\sqrt{N\mathfrak p}\rfloor
\ge2$.

Every split good prime of $k$ has two primes of $F$ of the same norm
and the same Frobenius trace. A prime of $F$ above an inert prime has
norm $(N\mathfrak p)^2$; consequently the number of such primes of
norm at most $x$ is $O_k(\sqrt x)$, even by the elementary bound for
prime ideals above rational primes. Ramified primes and primes where
good reduction changes under base extension form a fixed finite set.
It follows that
\[
 M^\pm_{E/F}(x)=2M^\pm_{E/k}(x)+O_{E,k}(\sqrt x).
\]
Apply the already proved connected-CM case over $F$ and divide by
two. This proves both formulas for $d=2$, with the same stated
error scale.

\paragraph{Literature comparison, geometry, and proof location.}
Bucur--Fit\'e--Kedlaya Proposition4.9 and Corollary4.10, published
pp.1741--1742, give the order $x^{3/4}/\log x$ and the quadratic
base-change mechanism under this same Hecke-power input.
James--Pollack, \emph{Extremal primes for elliptic curves with
complex multiplication} (2017), Theorem1 (author p.2), states the
$2/(3\pi)$ asymptotic over $\mathbb Q$ unconditionally; its separate
erratum replaces the narrow-sector analytic input. That published
result and its erratum are comparison references, not inputs to
the proof above. This run does not certify the erratum's sector
estimate independently.

The arithmetic content is the number of closed fibres for which the
trace on $H^1$ reaches the integral Hasse endpoint. The CM torus and
its rank-one Hecke characters supply all the analytic hypotheses;
there is no implication here for extremal traces of a non-CM curve.
Work/run: \texttt{GRH\_NIGHT/number\_theory},
\texttt{cm\_prime\_angles.tex}, NT-E08; complete deduction from NT-E06,
with source inputs S-NT-BFK and S-NT-ET and comparison S-NT-JP.
